\documentclass[11pt,a4paper]{article}
\usepackage[utf8]{inputenc}
\usepackage[T2A]{fontenc}
\usepackage[english,russian]{babel}
\usepackage{amsmath}
\usepackage{amssymb}
\usepackage[margin=1in]{geometry}
\usepackage{hyperref}
\usepackage{listings}
\usepackage{xcolor}
\usepackage{booktabs}
\usepackage{url}
% Моноширинная команда для длинных идентификаторов/путей, которая умеет
% переноситься на . _ / (обычный \texttt их не ломает и вылезает за поля).
\makeatletter
\DeclareUrlCommand\pathtt{\urlstyle{tt}}
\g@addto@macro\UrlBreaks{\do\.\do\_\do\/\do\-}
\makeatother

\hypersetup{colorlinks=true, linkcolor=blue!50!black,
            urlcolor=blue!50!black, citecolor=blue!50!black}

% OCaml listings style
\lstdefinelanguage{OCaml}{
  keywords={let,in,fun,match,with,type,module,sig,struct,
            and,or,not,if,then,else,begin,end,val,mutable,
            open,include,functor,of,rec,when},
  keywordstyle=\color[rgb]{0.0,0.4,0.7}\bfseries,
  commentstyle=\color[rgb]{0.4,0.4,0.4}\itshape,
  stringstyle=\color[rgb]{0.2,0.6,0.2},
  morecomment=[s]{(*}{*)},
  morestring=[b]",
  sensitive=true
}
\lstset{
  language=OCaml,
  basicstyle=\ttfamily\small,
  frame=single,
  framerule=0.4pt,
  rulecolor=\color[rgb]{0.8,0.8,0.8},
  backgroundcolor=\color[rgb]{0.97,0.97,0.97},
  xleftmargin=1em,
  xrightmargin=1em,
  aboveskip=0.8em,
  belowskip=0.8em,
  breaklines=true,
  showstringspaces=false,
  tabsize=2,
  extendedchars=true,
  literate=%
    {а}{{\cyra}}1 {б}{{\cyrb}}1 {в}{{\cyrv}}1 {г}{{\cyrg}}1
    {д}{{\cyrd}}1 {е}{{\cyre}}1 {ж}{{\cyrzh}}1 {з}{{\cyrz}}1
    {и}{{\cyri}}1 {й}{{\cyrishrt}}1 {к}{{\cyrk}}1 {л}{{\cyrl}}1
    {м}{{\cyrm}}1 {н}{{\cyrn}}1 {о}{{\cyro}}1 {п}{{\cyrp}}1
    {р}{{\cyrr}}1 {с}{{\cyrs}}1 {т}{{\cyrt}}1 {у}{{\cyru}}1
    {ф}{{\cyrf}}1 {х}{{\cyrh}}1 {ц}{{\cyrc}}1 {ч}{{\cyrch}}1
    {ш}{{\cyrsh}}1 {щ}{{\cyrshch}}1 {ъ}{{\cyrhrdsn}}1 {ы}{{\cyrery}}1
    {ь}{{\cyrsftsn}}1 {э}{{\cyrerev}}1 {ю}{{\cyryu}}1 {я}{{\cyrya}}1
    {А}{{\CYRA}}1 {Б}{{\CYRB}}1 {В}{{\CYRV}}1 {Г}{{\CYRG}}1
    {Д}{{\CYRD}}1 {Е}{{\CYRE}}1 {Ж}{{\CYRZH}}1 {З}{{\CYRZ}}1
    {И}{{\CYRI}}1 {Й}{{\CYRISHRT}}1 {К}{{\CYRK}}1 {Л}{{\CYRL}}1
    {М}{{\CYRM}}1 {Н}{{\CYRN}}1 {О}{{\CYRO}}1 {П}{{\CYRP}}1
    {Р}{{\CYRR}}1 {С}{{\CYRS}}1 {Т}{{\CYRT}}1 {У}{{\CYRU}}1
    {Ф}{{\CYRF}}1 {Х}{{\CYRH}}1 {Ц}{{\CYRC}}1 {Ч}{{\CYRCH}}1
    {Ш}{{\CYRSH}}1 {Щ}{{\CYRSHCH}}1 {Ъ}{{\CYRHRDSN}}1 {Ы}{{\CYRERY}}1
    {Ь}{{\CYRSFTSN}}1 {Э}{{\CYREREV}}1 {Ю}{{\CYRYU}}1 {Я}{{\CYRYA}}1
    {ё}{{\cyryo}}1 {Ё}{{\CYRYO}}1
}

\title{miniFlink: руководство для начала работы\\
\large Сквозной разбор от потока-функции до durable-пайплайна}
\author{Проект miniFlink}
\date{}

\begin{document}
\maketitle

\begin{abstract}
Это пошаговое руководство для тех, кто впервые видит miniFlink. Мы
строим один пайплайн от начала до конца на сквозном примере ---
телеметрии температурных датчиков --- и на каждом шаге добавляем ровно
одно понятие: поток, событие и время, watermark, ключ группировки,
окна, агрегацию, stateful-правила, дедупликацию и, наконец,
персистентность. Все листинги взяты из компилируемого кода. Прочитав
руководство сверху вниз, вы сможете читать и писать пайплайны miniFlink
и будете знать, куда смотреть дальше. Примеры по возрастанию сложности
живут в каталоге \texttt{examples/} (\texttt{ex01}--\texttt{ex10}); это
руководство --- сквозная нить, связывающая их.
\end{abstract}

\tableofcontents
\bigskip

%%% ──────────────────────────────────────────────────────────────────────
\section{Введение: зачем нужна потоковая обработка}

Этот раздел не про код. Он объясняет простыми словами, что такое
потоковая обработка, какую задачу она решает и почему для некоторых
задач она удобнее привычных подходов --- запросов к базе данных или
самостоятельно написанной программы-обработчика. Если вы никогда не
слышали этих слов --- ничего страшного, начнём с нуля.

\subsection{Два способа работать с данными}

Представьте, что вы следите за своим банковским счётом. Узнать, что с
деньгами что-то не так, можно двумя способами.

\textbf{Способ первый.} В конце месяца приходит выписка. Вы садитесь и
просматриваете все операции разом. Картина полная и аккуратная: все
траты, итоги, сравнение с прошлым месяцем. Но если в середине месяца
кто-то украл карту и потратил деньги, вы узнаёте об этом только сейчас
--- спустя недели, когда деньги уже ушли.

\textbf{Способ второй.} В ту самую секунду, когда происходит
подозрительная операция, вам приходит сообщение: <<с вашей карты
пытаются снять крупную сумму в другом городе>>. Вы реагируете
немедленно, пока ничего ещё не случилось.

Это не <<хороший>> и <<плохой>> способы. Это два разных инструмента для
двух разных задач. Первый --- \emph{пакетная} обработка (по-английски
batch): данные сначала накапливаются, а потом обрабатываются все сразу.
Второй --- \emph{потоковая} обработка: каждое событие обрабатывается в
тот момент, когда оно происходит, не дожидаясь остальных.

\subsection{Что такое <<поток>>}

Слово \emph{поток} здесь стоит понимать буквально --- как поток воды в
реке. У реки нет <<конца>>, который можно дождаться: вода течёт
постоянно. Точно так же во многих системах данные приходят
непрерывно: датчик каждую секунду присылает температуру, касса
пробивает чек за чеком, пользователи кликают по сайту круглые сутки.
Каждую такую <<каплю>> --- одно показание, один чек, один клик --- мы
будем называть \emph{событием}.

Пакетный подход говорит: <<накопим события за час (или за день), потом
обработаем>>. Потоковый подход говорит: <<обработаем каждое событие, как
только оно появилось, и сразу обновим результат>>. Разница --- как
между <<посмотреть утром вчерашнюю запись с камеры>> и <<охранник
смотрит на мониторы прямо сейчас и реагирует на ходу>>.

\subsection{Почему не просто запрос к базе раз в час?}

Самый привычный способ работать с данными --- сложить их в базу и
писать запросы (обычно на языке SQL): <<покажи среднюю температуру за
последний час>>, <<посчитай продажи за день>>. Для отчётов это
прекрасно. Но если задача --- \emph{реагировать вовремя}, у такого
подхода есть три слабых места.

\textbf{Задержка.} Если запрос запускается раз в час, то о проблеме вы
узнаёте в среднем через полчаса после того, как она случилась. Для
отчёта это неважно. Для тревоги о перегреве оборудования или утечке
газа --- полчаса может быть слишком поздно.

\textbf{Лишняя работа.} Каждый раз запрос пересчитывает всё заново ---
перебирает все накопившиеся данные, даже если изменилась всего пара
строк. Чем больше данных, тем тяжелее и дороже каждый прогон. Поток же
обрабатывает только новое событие и просто \emph{подправляет} уже
посчитанный результат.

\textbf{Самое тонкое --- время.} У каждого события есть два разных
момента: когда оно \emph{произошло на самом деле} и когда оно
\emph{дошло до вас}. В реальности данные опаздывают и приходят не по
порядку: датчик в шахте мог измерить температуру в 14:00, но из-за
плохой связи его показание добралось до сервера только в 14:05, уже
после более свежих данных. Если считать <<по времени получения>>, легко
сделать неверный вывод --- например, <<датчик молчал пять минут>>, хотя
он просто не мог достучаться. Потоковые системы умеют рассуждать о
времени \emph{события}, а не о времени его получения, и аккуратно
разбираться с опоздавшими данными. Сделать это правильно вручную ---
на удивление сложно.

\subsection{Почему не просто Kafka или <<напишу сам>>?}

Здесь часто звучит слово \emph{Kafka}. Полезно понимать, что это такое
и чем оно \emph{не} является.

Kafka --- это, если по-простому, очень надёжный конвейер (или почта)
для сообщений. Одна программа кладёт события на конвейер, другая ---
снимает их с другого конца. Kafka гарантирует, что сообщения не
потеряются и дойдут по порядку. Это отличный и важный инструмент --- но
он занимается \emph{перевозкой} событий, а не их \emph{осмыслением}.
Kafka сам по себе не знает, что такое <<среднее за минуту>>, не умеет
группировать события в окна времени, не отличает дубликат от нового
события и не подскажет, что делать с опоздавшим показанием.

Значит, всю эту логику кто-то должен написать в программе, которая
читает события с конвейера. И вот тут кроется главная трудность. Вещи,
которые на словах звучат просто, на практике полны подводных камней:

\begin{itemize}
\item считать <<среднее за последнюю минуту>> --- нужно решить, что
  такое <<минута>>, когда её закрывать и что делать с опоздавшими;
\item не посчитать одно и то же дважды, если событие случайно пришло
  по конвейеру два раза;
\item не сломаться, когда события приходят не по порядку;
\item не потерять накопленные итоги, если программа перезапустилась
  (упала и поднялась снова).
\end{itemize}

Каждый из этих пунктов --- источник тонких, незаметных ошибок, которые
проявляются редко и в самый неудачный момент. Написанные руками с нуля,
они почти наверняка где-то будут неверны.

\subsection{Что даёт потоковый движок}

Потоковый движок (Apache Flink --- самый известный; miniFlink ---
компактная реализация той же идеи) --- это набор готовых, проверенных
<<кирпичиков>> ровно для этих задач: окна времени, подсчёты и средние,
правила <<подними тревогу, когда значение перешло порог>>, защита от
дубликатов, аккуратная обработка опоздавших данных и сохранение
состояния при перезапуске. Вы не пишете эту механику заново --- вы
\emph{собираете} пайплайн из кирпичиков, а сложные и опасные детали уже
решены внутри и покрыты тестами.

Важно и то, \emph{как} это выглядит. Хорошо собранный потоковый
пайплайн читается почти как описание задачи на человеческом языке:
<<сгруппируй по датчику, посчитай среднее за минуту, подними тревогу при
перегреве>>. Не <<как именно перебрать данные>>, а <<что мы хотим
получить>>. Именно к этому --- чтобы код читался как спецификация ---
и стремится miniFlink. Дальше в руководстве вы увидите, что весь
пайплайн --- это короткая цепочка таких шагов.

\subsection{Когда пакетная обработка --- правильный выбор}

Будет честно сказать и обратное. Потоковая обработка --- не замена
всему. Если задача --- ночной отчёт, пересчёт зарплат раз в месяц,
аналитика по историческим данным или разовая выгрузка --- пакетный
подход и обычный SQL проще, понятнее и полностью достаточны. Реагировать
<<прямо сейчас>> там не нужно, а полная картина важнее скорости.

Потоковая обработка выигрывает там, где совпадают три вещи: данные
поступают \emph{непрерывно}, реагировать нужно \emph{быстро}, и важно
\emph{время самих событий} (а также корректность при опозданиях,
дубликатах и перезапусках). Мониторинг оборудования, обнаружение
мошенничества, живые дашборды, оповещения о безопасности --- классические
примеры. Система мониторинга шахты, на операторах которой построено это
руководство, --- ровно такой случай.

\subsection{Как читать дальше}

Дальше мы построим один такой пайплайн от начала до конца на сквозном
примере --- показаниях температурных датчиков. Каждый следующий раздел
добавляет ровно одно понятие из перечисленных выше: сначала
<<поток>>, потом <<событие и время>>, потом окна, тревоги, защиту от
дубликатов, обработку опоздавших данных и, наконец, сохранение
состояния. К концу вы сможете читать и писать такие пайплайны сами.
Технических терминов бояться не нужно --- каждый вводится по одному и
сразу на примере.

%%% ──────────────────────────────────────────────────────────────────────
\section{Прежде чем начать}

miniFlink --- это библиотека потоковой обработки на OCaml,
реализующая ядро семантики Apache Flink на одном узле: событийное
время, watermarks, окна, stateful-операторы, retract/update для
поздних данных и exactly-once. Никакого брокера, планировщика или
фонового цикла; весь пайплайн --- это композиция обычных функций,
которую можно читать как спецификацию.

Чтобы запускать примеры, нужен OCaml (4.14+ или 5.x) и Dune~3.0+.
Соберите проект командой \texttt{dune build} и прогоните тесты ---
\texttt{dune test}. Каждый листинг ниже предполагает, что в начале
файла стоит \texttt{open Miniflink}.

\paragraph{Сквозной пример.} На протяжении всего руководства мы
обрабатываем поток показаний температурных датчиков и хотим в итоге:
считать среднюю температуру по датчику в окне, поднимать тревогу при
перегреве и переживать перезапуск процесса без потери состояния.
Доменная модель проста:

\begin{lstlisting}
type reading = { sensor : string; celsius : float }
\end{lstlisting}

%%% ──────────────────────────────────────────────────────────────────────
\section{Шаг 1. Поток как функция}

Центральная абстракция --- \texttt{Stream.t}, и она устроена предельно
просто:

\begin{lstlisting}
type 'a t = unit -> 'a option
\end{lstlisting}

Поток --- это функция без аргументов, которая при каждом вызове
возвращает либо следующий элемент (\texttt{Some x}), либо
\texttt{None}, если элементы кончились. Это \emph{pull}-модель:
ничего не происходит, пока потребитель не «потянет» следующее
значение. Весь граф обработки --- это стопка таких функций, обёрнутых
одна вокруг другой.

\begin{lstlisting}
let nats : int Stream.t =
  let i = ref 0 in
  fun () -> let v = !i in incr i; if v < 5 then Some v else None

let doubled = Stream.map (fun x -> x * 2) nats
\end{lstlisting}

\texttt{Stream.map}, \texttt{Stream.filter} и \texttt{Stream.flat\_map}
--- это операторы преобразования: каждый принимает поток и возвращает
новый поток, не вычисляя ничего заранее. Композиция через оператор
\texttt{|>} читается естественно слева направо.

%%% ──────────────────────────────────────────────────────────────────────
\section{Шаг 2. Событие и время}

Голый поток значений ничего не знает о времени. В потоковой обработке
время --- первоклассное понятие, поэтому miniFlink оборачивает значения
в \texttt{Mf\_event.t} с четырьмя вариантами:

\begin{itemize}
\item \texttt{Data (value, ts)} --- значение с временной меткой
  (event-time);
\item \texttt{Watermark ts} --- отметка прогресса времени (см.\ шаг~3);
\item \texttt{Retract (value, ts)} --- отмена ранее выданного
  результата;
\item \texttt{Update \{old; new\_value; ts\}} --- атомарная замена
  старого результата новым.
\end{itemize}

Последние два понадобятся, когда мы дойдём до поздних данных; пока что
работаем с \texttt{Data}. Превратим наши показания в event-time поток
--- пусть каждое следующее чтение приходит на секунду позже:

\begin{lstlisting}
let event_stream : reading Mf_event.t Stream.t =
  readings
  |> List.mapi (fun i r -> Mf_event.data r (i * 1000))
  |> Stream.of_list
\end{lstlisting}

Временные метки здесь в миллисекундах; \texttt{Time} --- это просто
\texttt{int} мс с удобными конструкторами (\texttt{Time.seconds},
\texttt{Time.minutes}).

%%% ──────────────────────────────────────────────────────────────────────
\section{Шаг 3. Watermarks}

Событийное время порождает вопрос: когда можно считать, что «все
события до момента $t$ уже пришли»? В распределённой системе данные
опаздывают и приходят не по порядку. \emph{Watermark} --- это обещание
системы: «событий с меткой меньше $t$ больше не будет». Операторы окон
используют его, чтобы понять, когда окно можно закрывать и эмитить
результат.

\begin{lstlisting}
let with_wm =
  Mf_event.with_watermarks ~latency:(Time.seconds 1) event_stream
\end{lstlisting}

Параметр \texttt{\textasciitilde{}latency} задаёт допуск на опоздание:
watermark отстаёт от максимальной увиденной метки на эту величину, что
оставляет опоздавшим событиям шанс прийти. Чем больше latency, тем
устойчивее к опозданиям, но тем позже закрываются окна --- это прямой
компромисс между полнотой и задержкой.

%%% ──────────────────────────────────────────────────────────────────────
\section{Шаг 4. Ключ группировки через \texttt{KEYED}}

Почти все интересные операторы работают \emph{по ключу}: окно «на
датчик», агрегат «на пользователя». В большинстве фреймворков ключ
передаётся параметром в каждый оператор. miniFlink кодирует ключ
\emph{один раз} в типе через first-class-модуль \texttt{Keyed.S}:

\begin{lstlisting}
module By_sensor : Keyed.S with type t = reading = struct
  type t = reading
  let key r = r.sensor
end
\end{lstlisting}

Теперь \texttt{By\_sensor} можно передавать операторам как значение, и
им не нужен отдельный параметр \texttt{\textasciitilde{}key:}. Для
разовых случаев есть \texttt{Keyed.of\_fun (fun r -> r.sensor)}, не
требующий объявления модуля. Этот приём и делает пайплайны похожими на
спецификацию: ключ виден один раз, а не повторяется в каждом вызове.

%%% ──────────────────────────────────────────────────────────────────────
\section{Шаг 5. Окна и агрегация}

Окно режет бесконечный поток на конечные куски, к которым можно
применить агрегат. Простейшее --- \emph{тумблинг}: неперекрывающиеся
окна фиксированного размера. Посчитаем среднюю температуру по каждому
датчику в окнах по три секунды:

\begin{lstlisting}
let avg_per_sensor =
  event_stream
  |> Pipe.window_agg (module By_sensor)
       (Pipe.tumbling (Time.seconds 3))
       (Agg.mean (fun r -> r.celsius))
\end{lstlisting}

\texttt{Pipe.window\_agg} принимает модуль ключа, спецификацию окна
(\texttt{tumbling}, \texttt{sliding}, и т.\,д.) и \emph{агрегат}.
Агрегаты --- это композируемая библиотека: \texttt{count},
\texttt{sum}, \texttt{mean}, \texttt{count\_if}, а также комбинаторы
вроде \texttt{both} (считать два агрегата сразу) и \texttt{contramap}
(адаптировать входной тип). На выходе --- поток
\texttt{(string * 'r)}: ключ и результат агрегата за закрытое окно.
Результат появляется, когда watermark пересекает границу окна --- вот
зачем был нужен шаг~3.

%%% ──────────────────────────────────────────────────────────────────────
\section{Шаг 6. Stateful-правила: триггеры}

Окна отвечают на «сколько/какое среднее», но для системы оповещения
нужен \emph{stateful}-автомат: поднять тревогу, когда значение
пересекает порог, и снять её, когда вернулось в норму, не дёргая
оператора на каждом событии. Это \texttt{Trigger}:

\begin{lstlisting}
let alert_spec =
  Trigger.create
    ~name:"overheat"
    ~condition:(Trigger.greater_than 90.0)
    ~produce_alert:(fun ~key ~value ~ts:_ ->
      Printf.sprintf "ALERT %s=%.0f" key value)
    ~produce_recovery:(fun ~key ~ts:_ ->
      Printf.sprintf "OK %s" key)
    ()

let alerts =
  event_stream
  |> Stream.map (Mf_event.map_value (fun r -> (r.sensor, r.celsius)))
  |> Trigger.of_stream alert_spec
\end{lstlisting}

Триггер --- это конечный автомат на ключ. Когда температура датчика
превышает 90, он эмитит \texttt{Data} с алертом; когда падает обратно
--- \texttt{Retract} прежнего алерта и \texttt{Data} с recovery. Можно
задать debounce (\texttt{\textasciitilde{}problem\_for}), гистерезис
(\texttt{greater\_than\_with\_hysteresis}) и произвольные условия
(\texttt{custom}). Важное свойство, проверенное property-тестами:
активных алертов на ключ не больше одного, а каждый recovery парен
своему алерту --- автомат не «залипает».

%%% ──────────────────────────────────────────────────────────────────────
\section{Шаг 7. Дедупликация поздних дублей}

Семантика at-least-once у источников означает дубли: один пакет может
прийти дважды. \texttt{Pipe.dedup} подавляет повторы с одним ключом в
пределах окна \texttt{cooldown}:

\begin{lstlisting}
module By_pair : Keyed.S with type t = string * float = struct
  type t = string * float
  let key (k, _) = k
end

let deduped =
  event_stream
  |> Stream.map (Mf_event.map_value (fun r -> (r.sensor, r.celsius)))
  |> Pipe.dedup (module By_pair)
       ~rule:(fun (k, _) -> k)
       ~cooldown:(Time.seconds 5)
\end{lstlisting}

Свойство, зафиксированное property-тестом: между двумя прошедшими
событиями одного ключа разрыв по времени не меньше \texttt{cooldown},
и каждое прошедшее событие действительно было во входе.

%%% ──────────────────────────────────────────────────────────────────────
\section{Ещё операторы для реальных систем}

Шаги 1--7 --- это костяк. Но в настоящей системе мониторинга (как
minePASS, на котором основан сквозной пример) почти всегда нужны ещё
четыре приёма: дополнить событие справочными данными, сгруппировать
всплеск активности, объединить несколько потоков и заметить, что
источник \emph{замолчал}. Все они --- такие же операторы, как уже
знакомые.

\subsection{Обогащение из справочника: \texttt{enrich}}

Событие часто несёт только идентификатор --- например, номер датчика
--- а человеку нужно человекочитаемое: где этот датчик стоит. Такие
справочные данные («датчик A --- штрек 3») живут в \emph{таблице}, и
\texttt{Pipe.enrich} подмешивает их в проходящие события по ключу:

\begin{lstlisting}
let dict = Table.of_list [ ("A", "Штрек 3"); ("B", "Ствол") ]

let enriched =
  event_stream
  |> Pipe.enrich (module By_sensor)
       ~from:dict
       ~merge:(fun r loc -> match loc with
         | Some location -> { r with location }
         | None -> r)
\end{lstlisting}

\texttt{\textasciitilde{}from} --- таблица-справочник, \texttt{\textasciitilde{}merge}
--- как вписать найденное значение (или его отсутствие) в событие.
Таблицу можно собрать из списка (\texttt{Table.of\_list}), из
\texttt{Hashtbl} или \emph{из другого потока} (\texttt{Table.build}) ---
тогда справочник обновляется на лету по мере прихода новых данных.

\subsection{Окна по паузам: \texttt{session\_window}}

Тумблинг-окна из шага~5 режут время на равные куски. Но иногда
естественная единица --- не <<минута>>, а \emph{всплеск активности}:
группа событий, идущих подряд, отделённая от следующей группы паузой.
Например, серия показаний, пока техника работает, затем затишье, затем
новая серия. \texttt{session\_window} закрывает окно, когда после
события прошло больше \texttt{\textasciitilde{}gap} тишины:

\begin{lstlisting}
let sessions =
  event_stream
  |> Pipe.session_window (module By_sensor) ~gap:(Time.seconds 5)
\end{lstlisting}

На выходе --- \texttt{(ключ, список событий сессии)}. Если два события
одного датчика разделены паузой больше \texttt{gap}, они попадут в
разные сессии; если меньше --- в одну. Размер окна, таким образом,
определяют сами данные, а не часы.

\subsection{Объединение потоков: \texttt{keyed\_join}}

Часто про один объект приходят разные данные из разных источников:
температура из одного потока, влажность --- из другого, вибрация ---
из третьего. \texttt{keyed\_join} сводит несколько потоков в один,
сопоставляя события по ключу:

\begin{lstlisting}
let joined =
  Pipe.keyed_join (module By_sensor) [ temp_stream; humidity_stream ]
\end{lstlisting}

На каждое новое событие из любого входного потока на выходе появляется
\texttt{(ключ, список последних значений из каждого потока)} --- где
ещё не пришедшие значения отмечены как <<пусто>>. Так в одном месте
собирается полная свежая картина по объекту, даже если источники
обновляются вразнобой.

\subsection{Возраст молчания: \texttt{silence\_age}}

Для безопасности часто опаснее не <<плохое значение>>, а \emph{тишина}:
датчик, который перестал присылать данные, может означать обрыв связи
или поломку. \texttt{Item.silence\_age} следит за каждым ключом и
сообщает, сколько времени прошло с его последнего события:

\begin{lstlisting}
let silence =
  event_stream
  |> Item.silence_age ~by:(fun r -> r.sensor) ~tick:(Time.seconds 10)
\end{lstlisting}

На каждое реальное событие он эмитит \texttt{(ключ, 0)} --- <<возраст
обнулился>>, --- а затем каждые \texttt{\textasciitilde{}tick}
миллисекунд тишины --- \texttt{(ключ, возраст)} с растущим числом. Этот
поток дальше можно подать в триггер из шага~6: <<подними тревогу, если
датчик молчит дольше минуты>>. Получается оповещение об \emph{отсутствии}
данных --- то, что batch-запрос увидел бы в лучшем случае с большим
опозданием.

%%% ──────────────────────────────────────────────────────────────────────
\section{Порог из данных, а не из кода}

В шаге~6 порог тревоги был \emph{вшит} прямо в код:

\begin{lstlisting}
~condition:(Trigger.greater_than 90.0)
\end{lstlisting}

Число 90 написано раз и навсегда. Но в реальной системе пороги меняются:
для разных зон они разные, оператор хочет подкручивать их на ходу, не
перезапуская программу. Хочется, чтобы порог был \emph{настройкой}, а не
строчкой в коде.

Идея простая: пусть порог станет таким же \emph{данным}, как и само
показание. Тогда условие будет сравнивать два числа, которые пришли
вместе с событием. Заведём для этого тип, который несёт и замер, и
действующий порог:

\begin{lstlisting}
type checked = { zone : string; celsius : float; max_celsius : float }
\end{lstlisting}

И условие триггера теперь смотрит \emph{внутрь} значения, а не на
зашитую константу --- для этого есть \texttt{Trigger.custom}, который
принимает два произвольных предиката:

\begin{lstlisting}
let condition =
  Trigger.custom
    ~problem:(fun c -> c.celsius > c.max_celsius)
    ~recovery:(fun c -> c.celsius <= c.max_celsius)
\end{lstlisting}

Всё остальное в триггере не меняется: он по-прежнему ведёт отдельный
автомат на зону, держит не больше одной активной тревоги и парит
recovery к тревоге. Изменилось только одно --- порог переехал из кода в
данные. Остаётся вопрос: \emph{откуда} в событии берётся поле
\texttt{max\_celsius}? Рассмотрим два способа, от простого к
правильному.

\subsection{Способ 1: обогащение из таблицы порогов}

Первый способ опирается на уже знакомый \texttt{enrich} (см. предыдущий
раздел). Пороги хранятся в таблице <<зона $\to$ порог>>, а
\texttt{enrich} подмешивает текущий порог в каждое проходящее показание:

\begin{lstlisting}
let thresholds = Table.of_list [ ("A", 90.); ("B", 80.) ]

let alerts =
  readings
  (* пока порог неизвестен, кладём заведомо безопасное значение *)
  |> Stream.map (Mf_event.map_value (fun r ->
       { zone = r.zone; celsius = r.celsius; max_celsius = infinity }))
  |> Pipe.enrich (module Checked_by_zone)
       ~from:thresholds
       ~merge:(fun c m -> match m with
         | Some max_celsius -> { c with max_celsius }
         | None -> c)
  |> Stream.map (Mf_event.map_value (fun c -> (c.zone, c)))
  |> Trigger.of_stream spec
\end{lstlisting}

Таблицу можно собрать не из списка, а \emph{из потока} порогов
(\texttt{Table.build}) --- тогда она будет обновляться сама, как только
оператор пришлёт новое значение. Это и есть <<динамическая таблица
уставок>>.

\paragraph{На что обратить внимание.} Два момента, важных именно для
тревог.

Во-первых, случай \texttt{None} в \texttt{\textasciitilde{}merge} --- это
ситуация <<для зоны порог ещё не задан>>. Что тут правильно, решаете вы:
в примере выше мы подставляем \texttt{infinity} (без порога тревоги не
будет --- безопасный по умолчанию выбор для <<больше чем>>), но в иной
системе разумнее, наоборот, поднять тревогу <<зона без настройки>>.
Молча проигнорировать этот случай --- плохая идея.

Во-вторых, у этого способа есть тонкость со временем. \texttt{enrich}
берёт \emph{то значение, что лежит в таблице прямо сейчас}. Если порог и
показание приходят вразнобой, показание может проверяться против ещё не
обновлённого порога. Для отчётов это незаметно, для тревог --- иногда
важно. Поэтому есть второй способ.

\subsection{Способ 2: слияние потоков и состояние по зоне}

Более аккуратный подход --- относиться к порогам как к ещё одному
\emph{потоку событий}, наравне с показаниями, и свести оба потока в
один. Тогда у нас одна лента, упорядоченная по времени, где
перемешаны два вида событий: <<пришло новое показание>> и <<оператор
сменил порог>>.

Сначала опишем, что входом может быть одно из двух:

\begin{lstlisting}
type input = Reading of reading | Setpoint of setpoint
\end{lstlisting}

Сольём оба потока в один по времени с помощью \texttt{Mf\_event.union}:

\begin{lstlisting}
let merged =
  Mf_event.union
    (Stream.map (Mf_event.map_value (fun r -> Reading r)) readings)
    (Stream.map (Mf_event.map_value (fun s -> Setpoint s)) setpoints)
\end{lstlisting}

Теперь обработаем эту общую ленту \emph{stateful}-оператором
\texttt{process\_keyed} (см.\ под капотом он же стоит за триггером). На
каждую зону он держит маленькое состояние --- текущий порог --- и
обновляет его, когда приходит новый порог, либо сравнивает с ним, когда
приходит показание:

\begin{lstlisting}
let alerts =
  merged
  |> Pipe.process_keyed (module Input_by_zone)
       ~init:(fun () -> { max_c = None })   (* порога ещё нет *)
       ~on_event:(fun ctx zone st input ->
         match input with
         | Setpoint s ->
           st.max_c <- Some s.max_celsius   (* запомнили новый порог *)
         | Reading r ->
           match st.max_c with
           | Some m when r.celsius > m ->
             ctx.emit (alert zone r.celsius m)
           | _ -> ())
       ~on_timer:(fun _ _ _ _ _ -> ())
\end{lstlisting}

У этого способа есть важное преимущество, которого нет у первого.
Состояние пересматривается на \emph{любое} входное событие --- в том
числе на смену порога. Представьте: температура 85, порог 90 ---
тихо. Оператор снижает порог до 80. Те же 85 теперь должны стать
тревогой \emph{немедленно}, не дожидаясь нового замера. Первый способ
(обогащение) этого не заметит, пока не придёт следующее показание ---
триггер видит только поток показаний. А здесь смена порога --- это
полноценное событие на ленте, и оператор может тут же переоценить
ситуацию. Для тревог это часто решающий довод.

\subsection{Две тонкости, о которых легко забыть}

\paragraph{Поток порогов почти всегда молчит.} Показания идут потоком, а
порог оператор меняет раз в день. <<Молчащий>> поток порогов опасен для
событийного времени: пока по нему ничего не приходит, его watermark
(шаг~3) стоит на месте и тормозит весь объединённый пайплайн --- тревоги
по показаниям замирают в ожидании. Лекарство --- продвигать время
такого редкого потока по настенным часам через
\texttt{Mf\_event.with\_idle\_watermarks}: <<если из источника давно
ничего не было, считаем, что время всё равно идёт>>.

\paragraph{Гистерезис тоже можно вынести в данные.} В шаге~6 был
встроенный гистерезис (поднять тревогу при одном значении, снять при
другом, чтобы не дребезжало у самого порога), но с фиксированными
числами. Через \texttt{custom} его легко сделать настраиваемым: пусть
порог-событие несёт \emph{два} числа --- порог тревоги и порог снятия:

\begin{lstlisting}
~problem: (fun c -> c.celsius > c.alarm)
~recovery:(fun c -> c.celsius < c.clear)
\end{lstlisting}

Тогда и ширину гистерезиса можно настраивать на каждую зону и менять на
ходу --- то, чего фиксированная версия не умела.

\medskip
\noindent
Главное, что стоит унести из этого раздела: как только порог становится
данными, а не константой, открывается всё то же, что мы умеем делать с
данными --- хранить, обновлять потоком, объединять, переоценивать. Логика
тревоги при этом остаётся такой же короткой и читаемой.

%%% ──────────────────────────────────────────────────────────────────────
\section{Шаг 8. Поздние данные: retract и update}

Вернёмся к четырёхвариантному событию из шага~2. Что, если событие
опоздало и пришло уже после закрытия его окна? Наивный вариант ---
выбросить его (потеря данных) или выдать новый результат поверх
старого (дашборд «мигает» промежуточным состоянием). miniFlink делает
третье: атомарно \emph{корректирует} результат.

Когда в уже закрытое окно приходит позднее значение, оператор эмитит
\texttt{Update}\,$(o \to n)$ --- одно событие, заменяющее старый
результат $o$ новым $n$ без промежуточного состояния. Для потребителя
это выглядит как мгновенный переход, без «мигания». Фундаментальный
закон, который мы проверяем property-тестами для всех stateless-операторов:

\[
\mathtt{Update}\,(o \to n) \;\equiv\;
\mathtt{Retract}(o)\;\text{затем}\;\mathtt{Data}(n),
\]

то есть атомарный \texttt{Update} --- это в точности пара
«отозвать старое, выдать новое», только без промежуточного шага.
Писать обработку \texttt{Retract}/\texttt{Update} вручную обычно не
нужно: операторы (\texttt{map}, \texttt{filter}, \texttt{window\_agg},
\texttt{trigger}) делают это сами и согласованно.

%%% ──────────────────────────────────────────────────────────────────────
\section{Шаг 9. Персистентность без изменения пайплайна}

Финальное требование: пережить перезапуск процесса. Активная тревога
не должна исчезнуть из-за рестарта --- иначе диспетчер увидит мнимое
снятие.

Ключевая идея miniFlink --- \emph{ортогональная} персистентность: один
и тот же пайплайн работает с durability и без неё, режим выбирается
\emph{снаружи}, а не вшит в операторы. Сначала пишем пайплайн как
обычно:

\begin{lstlisting}
let pipeline src =
  src
  |> Stream.map (Mf_event.map_value (fun r -> (r.sensor, r.celsius)))
  |> Trigger.of_stream alert_spec
\end{lstlisting}

Без персистентности --- просто запускаем. С персистентностью ---
оборачиваем тот же вызов в durable-контекст:

\begin{lstlisting}
(* в памяти, ничего не сохраняется *)
let ephemeral () = pipeline event_stream

(* то же самое, но состояние durable *)
let durable () =
  let backend = Persistence_backend.of_memory (Hashtbl.create 16) in
  Runtime_context.with_context
    (Runtime_context.durable backend)
    (fun () -> pipeline event_stream)
\end{lstlisting}

Пайплайн не изменился ни на строку --- разница только в обёртке
\texttt{with\_context}. Сериализация тоже ортогональна: состояние
операторов сериализуется по умолчанию (через \texttt{Marshal}), так что
для crash-recovery не нужно писать ни одного кодека; для эволюции
схемы между деплоями можно задать реестр кодеков в контексте, опять же
не трогая пайплайн. В продакшене вместо in-memory backend подставляется
RocksDB; пайплайн при этом снова не меняется.

Это свойство --- что durable-прогон с «крашем» посередине даёт тот же
результат, что непрерывный --- мы проверяем property-тестом на сотнях
случайных потоков и точек краша.

%%% ──────────────────────────────────────────────────────────────────────
\section{Шаг 10. Масштабирование и exactly-once}

Всё выше --- однопоточный пайплайн. Для пропускной способности
\pathtt{Parallel.run_parallel_simple} шардирует поток по ключу на $N$
воркеров. Когда нужна не просто параллельность, а \emph{exactly-once}
с восстановлением после сбоя, используется
\pathtt{Checkpoint_parallel.run_exactly_once}: координатор
инжектирует barrier-маркеры (снимок в стиле Chandy--Lamport), каждый
воркер снапшотит своё состояние, а checkpoint фиксируется атомарно,
когда снимки собраны со всех воркеров. Транзакционный sink и
восстановление source по offset замыкают exactly-once на одном узле.

\begin{lstlisting}
Checkpoint_parallel.run_exactly_once
  ~workers:4 ~capacity:256 ~checkpoint_every:1000
  ~key_of ~make_state ~process ~source ~sink ~store ()
\end{lstlisting}

Это другой, более высокий уровень гарантий, чем per-operator
персистентность из шага~9: тот отвечает за восстановление состояния
оператора, этот --- за end-to-end exactly-once со стороны источника и
приёмника. Выбирайте по задаче: для большинства пайплайнов достаточно
ортогональной персистентности.

\paragraph{Замечание про параллелизм.} На OCaml~4 истинной
параллельности нет (runtime-lock), но конвейеризация через каналы
прячет задержки; измеренное ускорение --- $2.48\times$ на 4 воркерах.
На OCaml~5 с \texttt{Domain} появляется реальный параллелизм, и тот же
код пользуется им без изменений.

%%% ──────────────────────────────────────────────────────────────────────
\section{Что дальше}

Вы прошли путь от потока-функции до durable-пайплайна с тревогами.
Дальнейшие шаги:

\begin{itemize}
\item \textbf{Примеры.} Каталог \texttt{examples/} (\texttt{ex01}--%
  \texttt{ex10}) идёт по возрастанию сложности; \texttt{ex07} ---
  полноценный сервис мониторинга шахты с триангуляцией и
  retract-пайплайном газовых тревог.
\item \textbf{Справочник API.} Интерфейсные файлы \texttt{.mli}
  (особенно \texttt{pipe.mli} и \texttt{trigger.mli}) написаны как
  самостоятельные руководства с разбором краевых случаев.
\item \textbf{Глубже в семантику.} Статья \pathtt{miniflink_article}
  разбирает устройство движка; \pathtt{docs/orthogonal-persistence.md}
  --- модель персистентности; \pathtt{docs/atomic-update-event.md} ---
  семантику \texttt{Update}.
\item \textbf{Прикладной разбор.} Статья по \texttt{ex07}
  (\pathtt{docs/ex07_article}) показывает, как из этих операторов
  собирается реальная система безопасности.
\end{itemize}

Главная мысль, которую стоит унести: пайплайн miniFlink читается как
спецификация. Ключ объявлен один раз, время и поздние данные обработаны
операторами, а решения о персистентности и параллелизме приняты
снаружи и не засоряют логику. Это и есть та декларативность, ради
которой всё затевалось.

\end{document}
